\section{Conclusion}

For my thesis work, I propose a novel approach that uses diffusion models to generate dynamically near-admissible trajectories within the coupled, task-constrained state space of robot-environment interactions.
Through the implementation of a task-agnostic conditioning mechanism that encodes constraints via robot state space samples, I aim to demonstrate that diffusion models can enable the efficient generation of diverse and high-quality trajectory seeds that concurrently satisfy both kinematic and dynamic constraints, thereby substantially reducing computation demands on subsequent dynamically-constrained trajectory optimization.
By leveraging the strengths of both generative approaches and classical methods, my proposed project (in conjunction with my prior work) aims to expand the scope of robotic manipulation beyond geometric pick-and-place to include dynamic non-prehensile skills and motions governed by object inertial properties, impulsive contact, and motion-induced forces.
